\subsection{A sharp degree floor for this fixed frame}
The following negative result explains why further polynomial-center search
alone cannot lower the degree below the Round~18 value. It is a theorem
about a precisely fixed polynomial family, not a general Diophantine
complexity bound.

\begin{theorem}\label{thm:degree-floor}
Keep all the scalar registers defined above, the authentic six coefficients,
and the eleven independent coordinates fixed. For any nonzero fixed integer
$L$ and any integer polynomial $\Psi$ in the ten coordinates other than $n$,
\[
 \deg\mathcal N_{L,\Psi}\ge300.
\]
The center $\Psi_{300}$ with $L=32$ attains degree $300$. Thus $300$ is the
exact minimum in this fixed-frame, constant-weight, integer-center family.
The degree assertion does not require $L$ to be even; the logical
zero-set theorem above does.
\end{theorem}
\begin{proof}
All ports are constants of degree zero and $\beta\ne0$. A subscript zero
in this proof denotes the complete highest homogeneous component, not an
arbitrarily chosen leading monomial. Set
\[
 q_0=6\beta j,\quad X_0=12q_0w,\quad a_0=X_0v,
 \quad F_0=12wC-zq_0,
\]
\[
 k_0=R_0=-q_0^3F_0,\qquad \ell_0=X_0v^2.
\]
These are nonzero polynomials of degrees $1,2,3,2,5,4$, respectively.
Exact substitution in the literal registers gives
\begin{center}\small
\begin{tabular}{@{}lrl@{}}\toprule
Register & Degree & Highest homogeneous component\\\midrule
$A$ &3&$a_0$\\
$\Delta$ &6&$a_0^2$\\
$D_m$ &8&$a_0^2c^2$\\
$B$ &12&$ra_0^3c^2$\\
$U$ &11&$ra_0^2c^4$\\
$D_f$ &22&$ra_0^5c^6$\\
$Q,\delta$ &28&$q_*=ra_0^7c^6$\\
$V$ &37&$v_*=4r^2a_0^8c^{10}T$\\
$C_y$ &102&$q_*v_*^2$\\
$D_t$ &18&$\ell_0^2k_0^2$\\
$\mathcal B$ &18&$-k_0^2q_0^8$\\
$P$ &56&$-p_*^2$, where $p_*=q_0^4\ell_0^2k_0^3s$\\
$S$ &20&$s_*=128sra_0q_0^{15}$\\\bottomrule
\end{tabular}
\end{center}
In particular the highest form of $P$ is a negative square. To check this
sign directly, the terms $Y^2k^2(q^2-D_{\rm err})$ and
$-4Y^2q^2\eta^2$ in $\mathcal B$ have degree at most $16$, whereas
$-k^2q^8$ has degree $18$. Since
\[
 P=D_t+(\mathcal B-4)D_t^2s^2,
\]
its degree-$56$ part is exactly $-p_*^2$. This is the invariant that
prevents a degree-$300$ cancellation.

The carrier has degree $35$. Indeed its leading factors are
\[
 (z_c)_0=8a_0q_0^3T-4q_0^3X_0^2,\qquad
 (A_c)_0=((z_c)_0)^5.
\]
The two summands of $(z_c)_0$ cannot cancel because one contains $T$ and
the other does not.

Write $d=\deg\Psi$, treating the zero polynomial as degree $-\infty$.
If $d<54$, then $\deg Z\le73$ and $\deg p_a\le174$, whereas
$\deg p_b=124$ and $\deg p_c=176$. Thus the inner expression
$H_*=p_a+p_b-p_c$ has degree $176$. The cross term has degree at most
$298$, so $\deg\mathcal N=352$.

If $d>54$, then $\deg Z=20+d$ and $\deg p_a=68+2d>176$.
It dominates both other terms of $H_*$. Consequently
\[
 \deg\mathcal N=136+4d>352,
\]
because $\deg(4p_ap_b)=192+2d<136+4d$.

It remains to consider $d=54$. Let $\psi_0$ be the nonzero degree-$54$
part of $\Psi$. The leading forms of $p_a$ and $p_b$ are respectively
$q_*s_*^2\psi_0^2$ and $-q_*s_*^2p_*^2$, of degrees $176$ and $124$.
Therefore the degree-$300$ part of the negative cross term is
\begin{equation}\label{eq:positive-cross}
 (-4p_ap_b)_{300}=(2q_*s_*^2\psi_0p_*)^2\ne0.
\end{equation}
Put $h=\deg H_*$. If $h<150$, this cross term proves degree $300$.
If $h>150$, $H_*^2$ proves degree $2h>300$. If $h=150$, the degree-$300$
part is the sum of the square of the highest form of $H_*$ and the nonzero
square in \eqref{eq:positive-cross}. A sum of real polynomial squares with
one nonzero summand is not the zero polynomial. Thus in all cases
$\deg\mathcal N=\max(2h,300)\ge300$. Attainment is verified next.
\end{proof}
This theorem does not cover different carriers, changed auxiliary norms,
polynomial weights, coordinate changes in the other ten variables, or a
different encoding of refutation. It does not imply a global minimum even
among eleven-variable equations. Within its stated family, however, an
unsuccessful center search has a mathematical explanation rather than
merely an exhausted search budget.

\subsection{The attaining center and exact highest degree}
For completeness, the attained bound is checked at the actual coefficients,
not inferred from a syntactic degree estimate. Work first in the polynomial
ring with independent symbols $A,c,r,T,R$ and weights $(3,1,1,1,5)$.
Keep the exact identities
\[
 \Delta=A^2-1,\quad U=rc^2(1+\Delta c^2),\quad
 B=rA(1+\Delta c^2),\quad BD_f=B+\Delta A^2U^2.
\]
Define
\begin{align*}
 \lambda&=(8\Delta+4)U-rc^4,&\Psi_0&=\lambda V,\\
 h_*&=2U^3T(c^2-2),&b_*&=16A^3U^2T,\\
 e_*&=3rc^6U^2T.
\end{align*}
The attaining center is $\Psi_{300}=4(\Psi_0+h_*+b_*)+e_*$.
Its squared approximation error
$D_{300}=\Psi_{300}^2-32^2BD_fV^2$ follows from the exact identities
\begin{align*}
 D_{316}&=[8Ur(2c^2-c^4)+r^2c^8-64B]V^2,\\
 D_{308}&=D_{316}+2\Psi_0h_*+h_*^2,\\
 D_{304}&=D_{308}+2(\Psi_0+h_*)b_*+b_*^2,\\
 D_{300}&=16D_{304}+8(\Psi_0+h_*+b_*)e_*+e_*^2.
\end{align*}
These are literal polynomial identities obtained by squaring and using
$BD_f=B+\Delta A^2U^2$. Expand with
$U=r(c^4A^2+c^2-c^4)$ and
$V=D_f(c-R+4AUT)-c$. The highest weighted part of $D_{300}$ is
\[
 4096A^{16}c^{26}r^6T^2,
\]
of weight $82$; every remaining term has smaller weight. The supplied exact
coefficient expansion verifies the cancellation of every term of greater
weight, not only a numerical sample.

Set $E_{300}=\delta\Psi_{300}^2-32^2BD_fC_y$. Since $C_y=QV^2-1$ and
$\delta=Q-1$,
\[
 E_{300}=\delta D_{300}-32^2BD_f(V^2-1).
\]
The second summand has weight $108$; the first has highest part
$4096A^{23}c^{32}r^7T^2$ of weight $110$. Put $W_*=Sn-A_c$ of degree
$35$. The exact norm cut is
\[
 p_a+p_b-p_c=\delta W_*^2+2\delta S\Psi_{300}W_*
             +\delta S^2P+S^2E_{300}.
\]
The first three terms have degrees at most $98$, $137$ and $124$,
respectively, whereas the last has exact degree $150$. Its square thus has
degree $300$. A nonzero $T^4$ term in that degree, before expanding $a_0,q_0$,
is
\[
 2^{52}s^4r^{18}c^{64}T^4a_0^{50}q_0^{60}.
\]
The leading cross term has $T$-degree only two, so it cannot cancel this
monomial. This independently establishes attainment; it agrees with the
positive-square argument of Theorem~\ref{thm:degree-floor}.

\subsection{What this scalar iteration establishes}
The unchanged attaining control has eleven unrestricted integer variables,
exact degree 300 and 148 live binary operations. No scalar degree, dimension
or operation-count improvement is claimed in this iteration. The new result
is the exact degree minimum within the fixed constant-weight,
integer-center family of Theorem~\ref{thm:degree-floor}.

Going below that bound requires changing another part of the arithmetic
frame and proving its witness implications afresh. In particular, a
carrier congruence for a product cannot be cancelled using a coprimality
statement about a different pair of factors. The existing scale guard is
retained; no unproved guard deletion is used by an accepted endpoint.

The finite scalar checker binds the attaining graph to its authentic
coefficient payloads and verifies the complete highest homogeneous forms
used in the degree theorem. This exact algebra is distinct from the retained
unbounded Pell, interpolation and refutation calibration. Neither successful
replay nor the constrained degree minimum asserts consistency or any global
minimum among eleven-variable Diophantine equations.

